% GENERATED, edit the French unit or the translation table
\mxfront{terminale-e1}
\mxnewpage\phantomsection\label{sujet:1-1}\addcontentsline{toc}{chapter}{Test 1 · Practice 1}

\section*{Test 1 · Practice 1}

\mxbanner{PRACTICE PAPER · TRAINING}

\textbf{Duration: 2 h · Marked out of 20 · \mxTallied{ex-1-1} independent exercises.}\\Calculator allowed. Every answer must be justified; incomplete working will be taken into account. Give exact values; round approximations to $10^{-4}$, unless stated otherwise.

\smallskip\textbf{Name:}\enspace\dotfill\quad\textbf{Class:}\enspace\dotfill\quad\textbf{Date:}\enspace\dotfill

\mxtally{ex-1-1}
\subsection*{Exercise 1 · A discrete evolution \hfill 5 points}Consider the sequence $(u_n)$ defined by $u_0=1$ and, for every natural number $n$, $u_{n+1}=\frac{1}{2}u_n+4$.

\begin{enumerate}

\item Calculate $u_1$ and $u_2$. \hfill\textit{(0.5 pt)}

\item Prove by induction that, for all $n\in\N$, $u_n=8-7\left(\frac{1}{2}\right)^n$. \hfill\textit{(2 pt)}

\item Deduce that, for every natural number $n$, $1\le u_n<8$. \hfill\textit{(1 pt)}

\item Study the sign of $u_{n+1}-u_n$ and deduce the direction of variation of $(u_n)$. \hfill\textit{(1 pt)}

\item A student claims: ``There exists an integer $n$ such that $u_n=8$.'' Is this claim true? Justify. \hfill\textit{(0.5 pt)}

\end{enumerate}

\mxtally{ex-1-1}
\subsection*{Exercise 2 · Induction and bounds \hfill 5 points}Consider the sequence $(v_n)$ defined by $v_0=0$ and, for every natural number $n$, $v_{n+1}=\sqrt{72+v_n}$.

\begin{enumerate}

\item Calculate $v_1$ and $v_2$. Give the exact values. \hfill\textit{(0.5 pt)}

\item Prove by induction that all terms are defined and that, for all $n\in\N$, $0\le v_n\le9$. \hfill\textit{(1.5 pt)}

\item Prove by induction that, for all $n\in\N$, $v_{n+1}\ge v_n$. \hfill\textit{(1.5 pt)}

\item Prove that, for all $n\in\N$, $v_n<9$. \hfill\textit{(1 pt)}

\item Consider a property $P(n)$, defined for $n\in\N$. We have proved $P(1)$ and, for all $n\ge1$, the implication $P(n)\Rightarrow P(n+1)$. Do these two results allow us to conclude that $P(0)$ is true? Justify. \hfill\textit{(0.5 pt)}

\end{enumerate}

\mxnewpage\mxcontinued{Test 1 · Practice 1 · continued}

\mxtally{ex-1-1}
\subsection*{Exercise 3 · Vector calculations in a basis \hfill 5 points}In a basis $\pt{\vi;\vj;\vec k}$ of space, we define the vectors
\[\begin{aligned}\vec u&=\vi+\vec k,&\vec v&=\vj+\vec k,\\
\vec w&=4\vi+5\vj+9\vec k,&\vec z&=4\vi+5\vj+10\vec k.\end{aligned}\]

\begin{enumerate}

\item Justify that $\vec u$ and $\vec v$ are not collinear. \hfill\textit{(1 pt)}

\item Write $\vec w$ as a linear combination of $\vec u$ and $\vec v$. \hfill\textit{(0.5 pt)}

\item Are the vectors $\vec u$, $\vec v$ and $\vec z$ coplanar? Justify. \hfill\textit{(1.5 pt)}

\item Express the vector $5\vec u-4\vec v$ in the basis $\pt{\vi;\vj;\vec k}$. \hfill\textit{(1 pt)}

\item Choose a point $A$, then define $B$, $C$ and $D$ by $\vect{AB}=\vec u$, $\vect{AC}=\vec v$ and $\vect{AD}=\vec w$. Prove that $A$, $B$, $C$ and $D$ are coplanar. \hfill\textit{(1 pt)}

\end{enumerate}

\mxtally{ex-1-1}
\subsection*{Exercise 4 · Geometry of a parallelepiped \hfill 5 points}Consider a parallelepiped $ABCDEFGH$ that is not flat. The vertices $E$, $F$, $G$ and $H$ are the respective images of $A$, $B$, $C$ and $D$ under the translation by vector $\vect{AE}$.
Let $\vi=\vect{AB}$, $\vj=\vect{AD}$ and $\vec k=\vect{AE}$.
The points $I$, $J$ and $K$ are defined by
\[\vect{AI}=\frac1{11}\vi,\qquad\vect{AJ}=\frac1{11}\vj,\qquad\vect{AK}=\frac1{11}\vec k.\]

\begin{enumerate}

\item Express $\vect{AG}$ and $\vect{BH}$ in the basis $\pt{\vi;\vj;\vec k}$. \hfill\textit{(1 pt)}

\item Prove that the lines $(IJ)$ and $(BD)$ are parallel. \hfill\textit{(1 pt)}

\item Express $\vect{CG}$ in the basis, then determine whether $\vect{AB}$ and $\vect{CG}$ are collinear. Justify. \hfill\textit{(1 pt)}

\item Let $O$ be the centre of the parallelepiped. Prove that $O$ lies on the line $(AG)$. \hfill\textit{(1 pt)}

\item Prove that the points $I$, $J$ and $K$ are not collinear. \hfill\textit{(1 pt)}

\end{enumerate}

\ifmxteacher\mxnewpage

\section*{Solutions · Test 1 · Practice 1}\label{cor:1-1}

\subsection*{Exercise 1 · A discrete evolution}\begin{enumerate}

\item \textit{[0.5 pt]} The relation is applied successively for indices $0$ and $1$:\[u_1=\frac{1}{2}\times 1+4=\frac{9}{2},\qquad u_2=\frac{1}{2}\times \frac{9}{2}+4=\frac{25}{4}.\]

\item \textit{[2 pt]} For $n\in\N$, let $P(n)$ the property $u_n=8-7\left(\frac{1}{2}\right)^n$.
\par\smallskip\textbf{Base case.} For $n=0$, $8-7\times1=1=u_0$. So $P(0)$ is true.
\par\smallskip\textbf{Inductive step.} Fix an integer $n\ge0$ and assume $P(n)$ true. By the relation defining the sequence, then by the induction hypothesis,
\[\begin{aligned}u_{n+1}&=\frac{1}{2}u_n+4\\
&=\frac{1}{2}\left[8-7\left(\frac{1}{2}\right)^n\right]+4\\
&=4+4-7\left(\frac{1}{2}\right)^{n+1}\\
&=8-7\left(\frac{1}{2}\right)^{n+1}.
\end{aligned}\]
This is $P(n+1)$: the property is hereditary.
\par\smallskip\textbf{Conclusion.} The principle of induction gives $u_n=8-7\left(\frac{1}{2}\right)^n$ for all $n\in\N$.

\item \textit{[1 pt]} Since $0<\frac{1}{2}<1$, for all $n\in\N$ we have $0<(\frac{1}{2})^n\le1$. Multiplying by $-7<0$ reverses the inequalities:
\[-7\le-7(\frac{1}{2})^n<0.\]
Adding $8$ and using the previous formula, we obtain $1\le u_n<8$.

\item \textit{[1 pt]} For all $n\in\N$, the explicit formula gives
\[\begin{aligned}u_{n+1}-u_n&=-7(\frac{1}{2})^{n+1}+7(\frac{1}{2})^n\\&=7(1-\frac{1}{2})(\frac{1}{2})^n=\frac{7}{2}(\frac{1}{2})^n>0.\end{aligned}\]
The sequence is therefore strictly increasing.

\item \textit{[0.5 pt]} The claim is false. For each integer $n$, $8-u_n=7(\frac{1}{2})^n>0$. So no term is equal to $8$, whatever the integer $n$.

\end{enumerate}

\subsection*{Exercise 2 · Induction and bounds}\begin{enumerate}

\item \textit{[0.5 pt]} $v_1=\sqrt{72+0}=\sqrt{72}$ and $v_2=\sqrt{72+\sqrt{72}}$.

\item \textit{[1.5 pt]} Let $P(n)$: ``$v_n$ is defined and $0\le v_n\le 9$''.
\par\smallskip\textbf{Base case.} $v_0=0$ is defined and belongs to $[0,9]$; $P(0)$ is true.
\par\smallskip\textbf{Inductive step.} Fix $n\in\N$ and assume $P(n)$ true. Then
\[0\le 72\le 72+v_n\le 72+9=81.\]
The square root is therefore defined and, since this function is increasing on $[0,+\infty[$,
\[0\le v_{n+1}=\sqrt{72+v_n}\le\sqrt{81}=9.\]
Thus $P(n+1)$ is true.
\par\smallskip\textbf{Conclusion.} By induction, all terms are defined and belong to $[0,9]$.

\item \textit{[1.5 pt]} For $n\in\N$, let $R(n)$: ``$v_{n+1}\ge v_n$''.
\par\smallskip\textbf{Base case.} $v_1=\sqrt{72}>0=v_0$; $R(0)$ is true.
\par\smallskip\textbf{Inductive step.} Fix $n\ge0$ and assume $v_{n+1}\ge v_n$. Adding $72$ and then using the fact that the square root is increasing,
\[v_{n+2}=\sqrt{72+v_{n+1}}\ge\sqrt{72+v_n}=v_{n+1}.\]
This is $R(n+1)$.
\par\smallskip\textbf{Conclusion.} By induction, $v_{n+1}\ge v_n$ for all $n\ge0$: $(v_n)$ is increasing.

\item \textit{[1 pt]} Let $S(n)$: ``$v_n<9$''. At index $0$, $v_0=0<9$.
Fix $n\ge0$ and assume $S(n)$ true. The radicands being positive, the strict increase of the square root gives
\[v_{n+1}=\sqrt{72+v_n}<\sqrt{72+9}=\sqrt{81}=9.\]
So $S(n+1)$ is true. The base case and the inductive step establish by induction that $v_n<9$ for all $n\in\N$.

\item \textit{[0.5 pt]} No. For example $P(n):n\ge1$ is true for all integers $n\ge1$ but false at $0$.

\end{enumerate}

\subsection*{Exercise 3 · Vector calculations in a basis}\begin{enumerate}

\item \textit{[1 pt]} In an equality $\vec v=t\vec u$, the coefficient of $\vj $ would be $1=0$, a contradiction by uniqueness of coordinates in a basis.

\item \textit{[0.5 pt]} Expanding, $4\vec u+5\vec v=4(\vi +\vec k)+5(\vj +\vec k)=4\vi +5\vj +9\vec k=\vec w$.

\item \textit{[1.5 pt]} Since $\vec u,\vec v$ are not collinear, coplanarity would be equivalent to $\vec z=\lambda\vec u+\mu\vec v$. Uniqueness of the decomposition in the basis requires
\[\lambda=4,\qquad\mu=5,\qquad\lambda+\mu=10.\]
The first two equalities give $\lambda+\mu=9$, contradicting the third. The three vectors are therefore not coplanar.

\item \textit{[1 pt]} $5\vec u-4\vec v=5(\vi +\vec k)-4(\vj +\vec k)=5\vi -4\vj +\vec k$.

\item \textit{[1 pt]} $\vect{AB}$ and $\vect{AC}$ are non-collinear and $\vect{AD}$ is their linear combination: $D\in(ABC)$.

\end{enumerate}

\mxneed{24}

\subsection*{Exercise 4 · Geometry of a parallelepiped}\begin{enumerate}

\item \textit{[1 pt]} By Chasles' relation, $\vect{AG}=\vect{AB}+\vect{BC}+\vect{CG}=\vi +\vj +\vec k$. Similarly, $\vect{BH}=\vect{BA}+\vect{AD}+\vect{DH}=-\vi +\vj +\vec k$.

\item \textit{[1 pt]} By Chasles' relation, $\vect{IJ}=\vect{AJ}-\vect{AI}=\frac{\vj -\vi }{11}=\frac{\vect{BD}}{11}$. These vectors are non-zero and collinear: $(IJ)$ and $(BD)$ are parallel.

\item \textit{[1 pt]} We have $\vect{CG}=\vect{AE}=\vec k$ and $\vect{AB}=\vi $. These two basis vectors are not collinear: an equality $\vec k=t\vi $ would contradict the uniqueness of the decomposition in this basis.

\item \textit{[1 pt]} The diagonals of a parallelepiped bisect each other. Thus $\vect{AO}=\frac12\vect{AG}$, which proves $O\in(AG)$.

\item \textit{[1 pt]} The points satisfy $\vect{IJ}=\frac{-\vi +\vj }{11}$ and $\vect{IK}=\frac{-\vi +\vec k}{11}$. If $\vect{IK}=t\vect{IJ}$, the coefficient of $\vec k$ would give $\frac{1}{11}=0$, which is impossible. The vectors are not collinear, so $I,J,K$ are not collinear.

\end{enumerate}

\fi

\mxnewpage\phantomsection\label{sujet:1-2}\addcontentsline{toc}{chapter}{Test 1 · Practice 2}

\section*{Test 1 · Practice 2}

\mxbanner{PRACTICE PAPER · TRAINING}

\textbf{Duration: 2 h · Marked out of 20 · \mxTallied{ex-1-2} independent exercises.}\\Calculator allowed. Every answer must be justified; incomplete working will be taken into account. Give exact values; round approximations to $10^{-4}$, unless stated otherwise.

\smallskip\textbf{Name:}\enspace\dotfill\quad\textbf{Class:}\enspace\dotfill\quad\textbf{Date:}\enspace\dotfill

\mxtally{ex-1-2}
\subsection*{Exercise 1 · A discrete evolution \hfill 5 points}Consider the sequence $(u_n)$ defined by $u_0=2$ and, for every natural number $n$, $u_{n+1}=\frac{1}{5}u_n+\frac{36}{5}$.

\begin{enumerate}

\item Calculate $u_1$ and $u_2$. \hfill\textit{(0.5 pt)}

\item Prove by induction that, for all $n\in\N$, $u_n=9-7\left(\frac{1}{5}\right)^n$. \hfill\textit{(2 pt)}

\item Deduce that, for every natural number $n$, $2\le u_n<9$. \hfill\textit{(1 pt)}

\item Study the sign of $u_{n+1}-u_n$ and deduce the direction of variation of $(u_n)$. \hfill\textit{(1 pt)}

\item A student claims: ``There exists an integer $n$ such that $u_n=9$.'' Is this claim true? Justify. \hfill\textit{(0.5 pt)}

\end{enumerate}

\mxtally{ex-1-2}
\subsection*{Exercise 2 · Induction and bounds \hfill 5 points}Consider the sequence $(v_n)$ defined by $v_0=0$ and, for every natural number $n$, $v_{n+1}=\sqrt{90+v_n}$.

\begin{enumerate}

\item Calculate $v_1$ and $v_2$. Give the exact values. \hfill\textit{(0.5 pt)}

\item Prove by induction that all terms are defined and that, for all $n\in\N$, $0\le v_n\le10$. \hfill\textit{(1.5 pt)}

\item Prove by induction that, for all $n\in\N$, $v_{n+1}\ge v_n$. \hfill\textit{(1.5 pt)}

\item Prove that, for all $n\in\N$, $v_n<10$. \hfill\textit{(1 pt)}

\item Consider a property $P(n)$, defined for $n\in\N$. We have proved $P(1)$ and, for all $n\ge1$, the implication $P(n)\Rightarrow P(n+1)$. Do these two results allow us to conclude that $P(0)$ is true? Justify. \hfill\textit{(0.5 pt)}

\end{enumerate}

\mxnewpage\mxcontinued{Test 1 · Practice 2 · continued}

\mxtally{ex-1-2}
\subsection*{Exercise 3 · Vector calculations in a basis \hfill 5 points}In a basis $\pt{\vi;\vj;\vec k}$ of space, we define the vectors
\[\begin{aligned}\vec u&=\vi+\vec k,&\vec v&=\vj+\vec k,\\
\vec w&=3\vi+7\vj+10\vec k,&\vec z&=3\vi+7\vj+11\vec k.\end{aligned}\]

\begin{enumerate}

\item Justify that $\vec u$ and $\vec v$ are not collinear. \hfill\textit{(1 pt)}

\item Write $\vec w$ as a linear combination of $\vec u$ and $\vec v$. \hfill\textit{(0.5 pt)}

\item Are the vectors $\vec u$, $\vec v$ and $\vec z$ coplanar? Justify. \hfill\textit{(1.5 pt)}

\item Express the vector $7\vec u-3\vec v$ in the basis $\pt{\vi;\vj;\vec k}$. \hfill\textit{(1 pt)}

\item Choose a point $A$, then define $B$, $C$ and $D$ by $\vect{AB}=\vec u$, $\vect{AC}=\vec v$ and $\vect{AD}=\vec w$. Prove that $A$, $B$, $C$ and $D$ are coplanar. \hfill\textit{(1 pt)}

\end{enumerate}

\mxtally{ex-1-2}
\subsection*{Exercise 4 · Geometry of a parallelepiped \hfill 5 points}Consider a parallelepiped $ABCDEFGH$ that is not flat. The vertices $E$, $F$, $G$ and $H$ are the respective images of $A$, $B$, $C$ and $D$ under the translation by vector $\vect{AE}$.
Let $\vi=\vect{AB}$, $\vj=\vect{AD}$ and $\vec k=\vect{AE}$.
The points $I$, $J$ and $K$ are defined by
\[\vect{AI}=\frac1{18}\vi,\qquad\vect{AJ}=\frac1{18}\vj,\qquad\vect{AK}=\frac1{18}\vec k.\]

\begin{enumerate}

\item Express $\vect{AG}$ and $\vect{BH}$ in the basis $\pt{\vi;\vj;\vec k}$. \hfill\textit{(1 pt)}

\item Prove that the lines $(IJ)$ and $(BD)$ are parallel. \hfill\textit{(1 pt)}

\item Express $\vect{CG}$ in the basis, then determine whether $\vect{AB}$ and $\vect{CG}$ are collinear. Justify. \hfill\textit{(1 pt)}

\item Let $O$ be the centre of the parallelepiped. Prove that $O$ lies on the line $(AG)$. \hfill\textit{(1 pt)}

\item Prove that the points $I$, $J$ and $K$ are not collinear. \hfill\textit{(1 pt)}

\end{enumerate}

\ifmxteacher\mxnewpage

\section*{Solutions · Test 1 · Practice 2}\label{cor:1-2}

\subsection*{Exercise 1 · A discrete evolution}\begin{enumerate}

\item \textit{[0.5 pt]} The relation is applied successively for indices $0$ and $1$:\[u_1=\frac{1}{5}\times 2+\frac{36}{5}=\frac{38}{5},\qquad u_2=\frac{1}{5}\times \frac{38}{5}+\frac{36}{5}=\frac{218}{25}.\]

\item \textit{[2 pt]} For $n\in\N$, let $P(n)$ the property $u_n=9-7\left(\frac{1}{5}\right)^n$.
\par\smallskip\textbf{Base case.} For $n=0$, $9-7\times1=2=u_0$. So $P(0)$ is true.
\par\smallskip\textbf{Inductive step.} Fix an integer $n\ge0$ and assume $P(n)$ true. By the relation defining the sequence, then by the induction hypothesis,
\[\begin{aligned}u_{n+1}&=\frac{1}{5}u_n+\frac{36}{5}\\
&=\frac{1}{5}\left[9-7\left(\frac{1}{5}\right)^n\right]+\frac{36}{5}\\
&=\frac{9}{5}+\frac{36}{5}-7\left(\frac{1}{5}\right)^{n+1}\\
&=9-7\left(\frac{1}{5}\right)^{n+1}.
\end{aligned}\]
This is $P(n+1)$: the property is hereditary.
\par\smallskip\textbf{Conclusion.} The principle of induction gives $u_n=9-7\left(\frac{1}{5}\right)^n$ for all $n\in\N$.

\item \textit{[1 pt]} Since $0<\frac{1}{5}<1$, for all $n\in\N$ we have $0<(\frac{1}{5})^n\le1$. Multiplying by $-7<0$ reverses the inequalities:
\[-7\le-7(\frac{1}{5})^n<0.\]
Adding $9$ and using the previous formula, we obtain $2\le u_n<9$.

\item \textit{[1 pt]} For all $n\in\N$, the explicit formula gives
\[\begin{aligned}u_{n+1}-u_n&=-7(\frac{1}{5})^{n+1}+7(\frac{1}{5})^n\\&=7(1-\frac{1}{5})(\frac{1}{5})^n=\frac{28}{5}(\frac{1}{5})^n>0.\end{aligned}\]
The sequence is therefore strictly increasing.

\item \textit{[0.5 pt]} The claim is false. For each integer $n$, $9-u_n=7(\frac{1}{5})^n>0$. So no term is equal to $9$, whatever the integer $n$.

\end{enumerate}

\subsection*{Exercise 2 · Induction and bounds}\begin{enumerate}

\item \textit{[0.5 pt]} $v_1=\sqrt{90+0}=\sqrt{90}$ and $v_2=\sqrt{90+\sqrt{90}}$.

\item \textit{[1.5 pt]} Let $P(n)$: ``$v_n$ is defined and $0\le v_n\le 10$''.
\par\smallskip\textbf{Base case.} $v_0=0$ is defined and belongs to $[0,10]$; $P(0)$ is true.
\par\smallskip\textbf{Inductive step.} Fix $n\in\N$ and assume $P(n)$ true. Then
\[0\le 90\le 90+v_n\le 90+10=100.\]
The square root is therefore defined and, since this function is increasing on $[0,+\infty[$,
\[0\le v_{n+1}=\sqrt{90+v_n}\le\sqrt{100}=10.\]
Thus $P(n+1)$ is true.
\par\smallskip\textbf{Conclusion.} By induction, all terms are defined and belong to $[0,10]$.

\item \textit{[1.5 pt]} For $n\in\N$, let $R(n)$: ``$v_{n+1}\ge v_n$''.
\par\smallskip\textbf{Base case.} $v_1=\sqrt{90}>0=v_0$; $R(0)$ is true.
\par\smallskip\textbf{Inductive step.} Fix $n\ge0$ and assume $v_{n+1}\ge v_n$. Adding $90$ and then using the fact that the square root is increasing,
\[v_{n+2}=\sqrt{90+v_{n+1}}\ge\sqrt{90+v_n}=v_{n+1}.\]
This is $R(n+1)$.
\par\smallskip\textbf{Conclusion.} By induction, $v_{n+1}\ge v_n$ for all $n\ge0$: $(v_n)$ is increasing.

\item \textit{[1 pt]} Let $S(n)$: ``$v_n<10$''. At index $0$, $v_0=0<10$.
Fix $n\ge0$ and assume $S(n)$ true. The radicands being positive, the strict increase of the square root gives
\[v_{n+1}=\sqrt{90+v_n}<\sqrt{90+10}=\sqrt{100}=10.\]
So $S(n+1)$ is true. The base case and the inductive step establish by induction that $v_n<10$ for all $n\in\N$.

\item \textit{[0.5 pt]} No. For example $P(n):n\ge1$ is true for all integers $n\ge1$ but false at $0$.

\end{enumerate}

\subsection*{Exercise 3 · Vector calculations in a basis}\begin{enumerate}

\item \textit{[1 pt]} In an equality $\vec v=t\vec u$, the coefficient of $\vj $ would be $1=0$, a contradiction by uniqueness of coordinates in a basis.

\item \textit{[0.5 pt]} Expanding, $3\vec u+7\vec v=3(\vi +\vec k)+7(\vj +\vec k)=3\vi +7\vj +10\vec k=\vec w$.

\item \textit{[1.5 pt]} Since $\vec u,\vec v$ are not collinear, coplanarity would be equivalent to $\vec z=\lambda\vec u+\mu\vec v$. Uniqueness of the decomposition in the basis requires
\[\lambda=3,\qquad\mu=7,\qquad\lambda+\mu=11.\]
The first two equalities give $\lambda+\mu=10$, contradicting the third. The three vectors are therefore not coplanar.

\item \textit{[1 pt]} $7\vec u-3\vec v=7(\vi +\vec k)-3(\vj +\vec k)=7\vi -3\vj +4\vec k$.

\item \textit{[1 pt]} $\vect{AB}$ and $\vect{AC}$ are non-collinear and $\vect{AD}$ is their linear combination: $D\in(ABC)$.

\end{enumerate}

\mxneed{24}

\subsection*{Exercise 4 · Geometry of a parallelepiped}\begin{enumerate}

\item \textit{[1 pt]} By Chasles' relation, $\vect{AG}=\vect{AB}+\vect{BC}+\vect{CG}=\vi +\vj +\vec k$. Similarly, $\vect{BH}=\vect{BA}+\vect{AD}+\vect{DH}=-\vi +\vj +\vec k$.

\item \textit{[1 pt]} By Chasles' relation, $\vect{IJ}=\vect{AJ}-\vect{AI}=\frac{\vj -\vi }{18}=\frac{\vect{BD}}{18}$. These vectors are non-zero and collinear: $(IJ)$ and $(BD)$ are parallel.

\item \textit{[1 pt]} We have $\vect{CG}=\vect{AE}=\vec k$ and $\vect{AB}=\vi $. These two basis vectors are not collinear: an equality $\vec k=t\vi $ would contradict the uniqueness of the decomposition in this basis.

\item \textit{[1 pt]} The diagonals of a parallelepiped bisect each other. Thus $\vect{AO}=\frac12\vect{AG}$, which proves $O\in(AG)$.

\item \textit{[1 pt]} The points satisfy $\vect{IJ}=\frac{-\vi +\vj }{18}$ and $\vect{IK}=\frac{-\vi +\vec k}{18}$. If $\vect{IK}=t\vect{IJ}$, the coefficient of $\vec k$ would give $\frac{1}{18}=0$, which is impossible. The vectors are not collinear, so $I,J,K$ are not collinear.

\end{enumerate}

\fi

\ifmxteacher

\mxnewpage\phantomsection\label{sujet:1-3}\addcontentsline{toc}{chapter}{Test 1 · A}

\section*{Test 1 · A}

\mxbanner*{PAPER RESERVED FOR ASSESSMENT}

\textbf{Duration: 2 h · Marked out of 20 · \mxTallied{ex-1-3} independent exercises.}\\Calculator allowed. Every answer must be justified; incomplete working will be taken into account. Give exact values; round approximations to $10^{-4}$, unless stated otherwise.

\smallskip\textbf{Name:}\enspace\dotfill\quad\textbf{Class:}\enspace\dotfill\quad\textbf{Date:}\enspace\dotfill

\mxtally{ex-1-3}
\subsection*{Exercise 1 · A discrete evolution \hfill 5 points}Consider the sequence $(u_n)$ defined by $u_0=3$ and, for every natural number $n$, $u_{n+1}=\frac{1}{4}u_n+\frac{33}{4}$.

\begin{enumerate}

\item Calculate $u_1$ and $u_2$. \hfill\textit{(0.5 pt)}

\item Prove by induction that, for all $n\in\N$, $u_n=11-8\left(\frac{1}{4}\right)^n$. \hfill\textit{(2 pt)}

\item Deduce that, for every natural number $n$, $3\le u_n<11$. \hfill\textit{(1 pt)}

\item Study the sign of $u_{n+1}-u_n$ and deduce the direction of variation of $(u_n)$. \hfill\textit{(1 pt)}

\item A student claims: ``There exists an integer $n$ such that $u_n=11$.'' Is this claim true? Justify. \hfill\textit{(0.5 pt)}

\end{enumerate}

\mxtally{ex-1-3}
\subsection*{Exercise 2 · Induction and bounds \hfill 5 points}Consider the sequence $(v_n)$ defined by $v_0=0$ and, for every natural number $n$, $v_{n+1}=\sqrt{110+v_n}$.

\begin{enumerate}

\item Calculate $v_1$ and $v_2$. Give the exact values. \hfill\textit{(0.5 pt)}

\item Prove by induction that all terms are defined and that, for all $n\in\N$, $0\le v_n\le11$. \hfill\textit{(1.5 pt)}

\item Prove by induction that, for all $n\in\N$, $v_{n+1}\ge v_n$. \hfill\textit{(1.5 pt)}

\item Prove that, for all $n\in\N$, $v_n<11$. \hfill\textit{(1 pt)}

\item Consider a property $P(n)$, defined for $n\in\N$. We have proved $P(1)$ and, for all $n\ge1$, the implication $P(n)\Rightarrow P(n+1)$. Do these two results allow us to conclude that $P(0)$ is true? Justify. \hfill\textit{(0.5 pt)}

\end{enumerate}

\mxnewpage\mxcontinued{Test 1 · A · continued}

\mxtally{ex-1-3}
\subsection*{Exercise 3 · Vector calculations in a basis \hfill 5 points}In a basis $\pt{\vi;\vj;\vec k}$ of space, we define the vectors
\[\begin{aligned}\vec u&=\vi+\vec k,&\vec v&=\vj+\vec k,\\
\vec w&=2\vi+4\vj+6\vec k,&\vec z&=2\vi+4\vj+7\vec k.\end{aligned}\]

\begin{enumerate}

\item Justify that $\vec u$ and $\vec v$ are not collinear. \hfill\textit{(1 pt)}

\item Write $\vec w$ as a linear combination of $\vec u$ and $\vec v$. \hfill\textit{(0.5 pt)}

\item Are the vectors $\vec u$, $\vec v$ and $\vec z$ coplanar? Justify. \hfill\textit{(1.5 pt)}

\item Express the vector $4\vec u-2\vec v$ in the basis $\pt{\vi;\vj;\vec k}$. \hfill\textit{(1 pt)}

\item Choose a point $A$, then define $B$, $C$ and $D$ by $\vect{AB}=\vec u$, $\vect{AC}=\vec v$ and $\vect{AD}=\vec w$. Prove that $A$, $B$, $C$ and $D$ are coplanar. \hfill\textit{(1 pt)}

\end{enumerate}

\mxtally{ex-1-3}
\subsection*{Exercise 4 · Geometry of a parallelepiped \hfill 5 points}Consider a parallelepiped $ABCDEFGH$ that is not flat. The vertices $E$, $F$, $G$ and $H$ are the respective images of $A$, $B$, $C$ and $D$ under the translation by vector $\vect{AE}$.
Let $\vi=\vect{AB}$, $\vj=\vect{AD}$ and $\vec k=\vect{AE}$.
The points $I$, $J$ and $K$ are defined by
\[\vect{AI}=\frac1{8}\vi,\qquad\vect{AJ}=\frac1{8}\vj,\qquad\vect{AK}=\frac1{8}\vec k.\]

\begin{enumerate}

\item Express $\vect{AG}$ and $\vect{BH}$ in the basis $\pt{\vi;\vj;\vec k}$. \hfill\textit{(1 pt)}

\item Prove that the lines $(IJ)$ and $(BD)$ are parallel. \hfill\textit{(1 pt)}

\item Express $\vect{CG}$ in the basis, then determine whether $\vect{AB}$ and $\vect{CG}$ are collinear. Justify. \hfill\textit{(1 pt)}

\item Let $O$ be the centre of the parallelepiped. Prove that $O$ lies on the line $(AG)$. \hfill\textit{(1 pt)}

\item Prove that the points $I$, $J$ and $K$ are not collinear. \hfill\textit{(1 pt)}

\end{enumerate}

\ifmxteacher\mxnewpage

\section*{Solutions · Test 1 · A}\label{cor:1-3}

\subsection*{Exercise 1 · A discrete evolution}\begin{enumerate}

\item \textit{[0.5 pt]} The relation is applied successively for indices $0$ and $1$:\[u_1=\frac{1}{4}\times 3+\frac{33}{4}=9,\qquad u_2=\frac{1}{4}\times 9+\frac{33}{4}=\frac{21}{2}.\]

\item \textit{[2 pt]} For $n\in\N$, let $P(n)$ the property $u_n=11-8\left(\frac{1}{4}\right)^n$.
\par\smallskip\textbf{Base case.} For $n=0$, $11-8\times1=3=u_0$. So $P(0)$ is true.
\par\smallskip\textbf{Inductive step.} Fix an integer $n\ge0$ and assume $P(n)$ true. By the relation defining the sequence, then by the induction hypothesis,
\[\begin{aligned}u_{n+1}&=\frac{1}{4}u_n+\frac{33}{4}\\
&=\frac{1}{4}\left[11-8\left(\frac{1}{4}\right)^n\right]+\frac{33}{4}\\
&=\frac{11}{4}+\frac{33}{4}-8\left(\frac{1}{4}\right)^{n+1}\\
&=11-8\left(\frac{1}{4}\right)^{n+1}.
\end{aligned}\]
This is $P(n+1)$: the property is hereditary.
\par\smallskip\textbf{Conclusion.} The principle of induction gives $u_n=11-8\left(\frac{1}{4}\right)^n$ for all $n\in\N$.

\item \textit{[1 pt]} Since $0<\frac{1}{4}<1$, for all $n\in\N$ we have $0<(\frac{1}{4})^n\le1$. Multiplying by $-8<0$ reverses the inequalities:
\[-8\le-8(\frac{1}{4})^n<0.\]
Adding $11$ and using the previous formula, we obtain $3\le u_n<11$.

\item \textit{[1 pt]} For all $n\in\N$, the explicit formula gives
\[\begin{aligned}u_{n+1}-u_n&=-8(\frac{1}{4})^{n+1}+8(\frac{1}{4})^n\\&=8(1-\frac{1}{4})(\frac{1}{4})^n=6(\frac{1}{4})^n>0.\end{aligned}\]
The sequence is therefore strictly increasing.

\item \textit{[0.5 pt]} The claim is false. For each integer $n$, $11-u_n=8(\frac{1}{4})^n>0$. So no term is equal to $11$, whatever the integer $n$.

\end{enumerate}

\subsection*{Exercise 2 · Induction and bounds}\begin{enumerate}

\item \textit{[0.5 pt]} $v_1=\sqrt{110+0}=\sqrt{110}$ and $v_2=\sqrt{110+\sqrt{110}}$.

\item \textit{[1.5 pt]} Let $P(n)$: ``$v_n$ is defined and $0\le v_n\le 11$''.
\par\smallskip\textbf{Base case.} $v_0=0$ is defined and belongs to $[0,11]$; $P(0)$ is true.
\par\smallskip\textbf{Inductive step.} Fix $n\in\N$ and assume $P(n)$ true. Then
\[0\le 110\le 110+v_n\le 110+11=121.\]
The square root is therefore defined and, since this function is increasing on $[0,+\infty[$,
\[0\le v_{n+1}=\sqrt{110+v_n}\le\sqrt{121}=11.\]
Thus $P(n+1)$ is true.
\par\smallskip\textbf{Conclusion.} By induction, all terms are defined and belong to $[0,11]$.

\item \textit{[1.5 pt]} For $n\in\N$, let $R(n)$: ``$v_{n+1}\ge v_n$''.
\par\smallskip\textbf{Base case.} $v_1=\sqrt{110}>0=v_0$; $R(0)$ is true.
\par\smallskip\textbf{Inductive step.} Fix $n\ge0$ and assume $v_{n+1}\ge v_n$. Adding $110$ and then using the fact that the square root is increasing,
\[v_{n+2}=\sqrt{110+v_{n+1}}\ge\sqrt{110+v_n}=v_{n+1}.\]
This is $R(n+1)$.
\par\smallskip\textbf{Conclusion.} By induction, $v_{n+1}\ge v_n$ for all $n\ge0$: $(v_n)$ is increasing.

\item \textit{[1 pt]} Let $S(n)$: ``$v_n<11$''. At index $0$, $v_0=0<11$.
Fix $n\ge0$ and assume $S(n)$ true. The radicands being positive, the strict increase of the square root gives
\[v_{n+1}=\sqrt{110+v_n}<\sqrt{110+11}=\sqrt{121}=11.\]
So $S(n+1)$ is true. The base case and the inductive step establish by induction that $v_n<11$ for all $n\in\N$.

\item \textit{[0.5 pt]} No. For example $P(n):n\ge1$ is true for all integers $n\ge1$ but false at $0$.

\end{enumerate}

\subsection*{Exercise 3 · Vector calculations in a basis}\begin{enumerate}

\item \textit{[1 pt]} In an equality $\vec v=t\vec u$, the coefficient of $\vj $ would be $1=0$, a contradiction by uniqueness of coordinates in a basis.

\item \textit{[0.5 pt]} Expanding, $2\vec u+4\vec v=2(\vi +\vec k)+4(\vj +\vec k)=2\vi +4\vj +6\vec k=\vec w$.

\item \textit{[1.5 pt]} Since $\vec u,\vec v$ are not collinear, coplanarity would be equivalent to $\vec z=\lambda\vec u+\mu\vec v$. Uniqueness of the decomposition in the basis requires
\[\lambda=2,\qquad\mu=4,\qquad\lambda+\mu=7.\]
The first two equalities give $\lambda+\mu=6$, contradicting the third. The three vectors are therefore not coplanar.

\item \textit{[1 pt]} $4\vec u-2\vec v=4(\vi +\vec k)-2(\vj +\vec k)=4\vi -2\vj +2\vec k$.

\item \textit{[1 pt]} $\vect{AB}$ and $\vect{AC}$ are non-collinear and $\vect{AD}$ is their linear combination: $D\in(ABC)$.

\end{enumerate}

\mxneed{24}

\subsection*{Exercise 4 · Geometry of a parallelepiped}\begin{enumerate}

\item \textit{[1 pt]} By Chasles' relation, $\vect{AG}=\vect{AB}+\vect{BC}+\vect{CG}=\vi +\vj +\vec k$. Similarly, $\vect{BH}=\vect{BA}+\vect{AD}+\vect{DH}=-\vi +\vj +\vec k$.

\item \textit{[1 pt]} By Chasles' relation, $\vect{IJ}=\vect{AJ}-\vect{AI}=\frac{\vj -\vi }{8}=\frac{\vect{BD}}{8}$. These vectors are non-zero and collinear: $(IJ)$ and $(BD)$ are parallel.

\item \textit{[1 pt]} We have $\vect{CG}=\vect{AE}=\vec k$ and $\vect{AB}=\vi $. These two basis vectors are not collinear: an equality $\vec k=t\vi $ would contradict the uniqueness of the decomposition in this basis.

\item \textit{[1 pt]} The diagonals of a parallelepiped bisect each other. Thus $\vect{AO}=\frac12\vect{AG}$, which proves $O\in(AG)$.

\item \textit{[1 pt]} The points satisfy $\vect{IJ}=\frac{-\vi +\vj }{8}$ and $\vect{IK}=\frac{-\vi +\vec k}{8}$. If $\vect{IK}=t\vect{IJ}$, the coefficient of $\vec k$ would give $\frac{1}{8}=0$, which is impossible. The vectors are not collinear, so $I,J,K$ are not collinear.

\end{enumerate}

\fi

\fi

\ifmxteacher

\mxnewpage\phantomsection\label{sujet:1-4}\addcontentsline{toc}{chapter}{Test 1 · B}

\section*{Test 1 · B}

\mxbanner*{PAPER RESERVED FOR ASSESSMENT}

\textbf{Duration: 2 h · Marked out of 20 · \mxTallied{ex-1-4} independent exercises.}\\Calculator allowed. Every answer must be justified; incomplete working will be taken into account. Give exact values; round approximations to $10^{-4}$, unless stated otherwise.

\smallskip\textbf{Name:}\enspace\dotfill\quad\textbf{Class:}\enspace\dotfill\quad\textbf{Date:}\enspace\dotfill

\mxtally{ex-1-4}
\subsection*{Exercise 1 · A discrete evolution \hfill 5 points}Consider the sequence $(u_n)$ defined by $u_0=1$ and, for every natural number $n$, $u_{n+1}=\frac{1}{3}u_n+\frac{16}{3}$.

\begin{enumerate}

\item Calculate $u_1$ and $u_2$. \hfill\textit{(0.5 pt)}

\item Prove by induction that, for all $n\in\N$, $u_n=8-7\left(\frac{1}{3}\right)^n$. \hfill\textit{(2 pt)}

\item Deduce that, for every natural number $n$, $1\le u_n<8$. \hfill\textit{(1 pt)}

\item Study the sign of $u_{n+1}-u_n$ and deduce the direction of variation of $(u_n)$. \hfill\textit{(1 pt)}

\item A student claims: ``There exists an integer $n$ such that $u_n=8$.'' Is this claim true? Justify. \hfill\textit{(0.5 pt)}

\end{enumerate}

\mxtally{ex-1-4}
\subsection*{Exercise 2 · Induction and bounds \hfill 5 points}Consider the sequence $(v_n)$ defined by $v_0=0$ and, for every natural number $n$, $v_{n+1}=\sqrt{42+v_n}$.

\begin{enumerate}

\item Calculate $v_1$ and $v_2$. Give the exact values. \hfill\textit{(0.5 pt)}

\item Prove by induction that all terms are defined and that, for all $n\in\N$, $0\le v_n\le7$. \hfill\textit{(1.5 pt)}

\item Prove by induction that, for all $n\in\N$, $v_{n+1}\ge v_n$. \hfill\textit{(1.5 pt)}

\item Prove that, for all $n\in\N$, $v_n<7$. \hfill\textit{(1 pt)}

\item Consider a property $P(n)$, defined for $n\in\N$. We have proved $P(1)$ and, for all $n\ge1$, the implication $P(n)\Rightarrow P(n+1)$. Do these two results allow us to conclude that $P(0)$ is true? Justify. \hfill\textit{(0.5 pt)}

\end{enumerate}

\mxnewpage\mxcontinued{Test 1 · B · continued}

\mxtally{ex-1-4}
\subsection*{Exercise 3 · Vector calculations in a basis \hfill 5 points}In a basis $\pt{\vi;\vj;\vec k}$ of space, we define the vectors
\[\begin{aligned}\vec u&=\vi+\vec k,&\vec v&=\vj+\vec k,\\
\vec w&=5\vi+5\vj+10\vec k,&\vec z&=5\vi+5\vj+11\vec k.\end{aligned}\]

\begin{enumerate}

\item Justify that $\vec u$ and $\vec v$ are not collinear. \hfill\textit{(1 pt)}

\item Write $\vec w$ as a linear combination of $\vec u$ and $\vec v$. \hfill\textit{(0.5 pt)}

\item Are the vectors $\vec u$, $\vec v$ and $\vec z$ coplanar? Justify. \hfill\textit{(1.5 pt)}

\item Express the vector $5\vec u-5\vec v$ in the basis $\pt{\vi;\vj;\vec k}$. \hfill\textit{(1 pt)}

\item Choose a point $A$, then define $B$, $C$ and $D$ by $\vect{AB}=\vec u$, $\vect{AC}=\vec v$ and $\vect{AD}=\vec w$. Prove that $A$, $B$, $C$ and $D$ are coplanar. \hfill\textit{(1 pt)}

\end{enumerate}

\mxtally{ex-1-4}
\subsection*{Exercise 4 · Geometry of a parallelepiped \hfill 5 points}Consider a parallelepiped $ABCDEFGH$ that is not flat. The vertices $E$, $F$, $G$ and $H$ are the respective images of $A$, $B$, $C$ and $D$ under the translation by vector $\vect{AE}$.
Let $\vi=\vect{AB}$, $\vj=\vect{AD}$ and $\vec k=\vect{AE}$.
The points $I$, $J$ and $K$ are defined by
\[\vect{AI}=\frac1{15}\vi,\qquad\vect{AJ}=\frac1{15}\vj,\qquad\vect{AK}=\frac1{15}\vec k.\]

\begin{enumerate}

\item Express $\vect{AG}$ and $\vect{BH}$ in the basis $\pt{\vi;\vj;\vec k}$. \hfill\textit{(1 pt)}

\item Prove that the lines $(IJ)$ and $(BD)$ are parallel. \hfill\textit{(1 pt)}

\item Express $\vect{CG}$ in the basis, then determine whether $\vect{AB}$ and $\vect{CG}$ are collinear. Justify. \hfill\textit{(1 pt)}

\item Let $O$ be the centre of the parallelepiped. Prove that $O$ lies on the line $(AG)$. \hfill\textit{(1 pt)}

\item Prove that the points $I$, $J$ and $K$ are not collinear. \hfill\textit{(1 pt)}

\end{enumerate}

\ifmxteacher\mxnewpage

\section*{Solutions · Test 1 · B}\label{cor:1-4}

\subsection*{Exercise 1 · A discrete evolution}\begin{enumerate}

\item \textit{[0.5 pt]} The relation is applied successively for indices $0$ and $1$:\[u_1=\frac{1}{3}\times 1+\frac{16}{3}=\frac{17}{3},\qquad u_2=\frac{1}{3}\times \frac{17}{3}+\frac{16}{3}=\frac{65}{9}.\]

\item \textit{[2 pt]} For $n\in\N$, let $P(n)$ the property $u_n=8-7\left(\frac{1}{3}\right)^n$.
\par\smallskip\textbf{Base case.} For $n=0$, $8-7\times1=1=u_0$. So $P(0)$ is true.
\par\smallskip\textbf{Inductive step.} Fix an integer $n\ge0$ and assume $P(n)$ true. By the relation defining the sequence, then by the induction hypothesis,
\[\begin{aligned}u_{n+1}&=\frac{1}{3}u_n+\frac{16}{3}\\
&=\frac{1}{3}\left[8-7\left(\frac{1}{3}\right)^n\right]+\frac{16}{3}\\
&=\frac{8}{3}+\frac{16}{3}-7\left(\frac{1}{3}\right)^{n+1}\\
&=8-7\left(\frac{1}{3}\right)^{n+1}.
\end{aligned}\]
This is $P(n+1)$: the property is hereditary.
\par\smallskip\textbf{Conclusion.} The principle of induction gives $u_n=8-7\left(\frac{1}{3}\right)^n$ for all $n\in\N$.

\item \textit{[1 pt]} Since $0<\frac{1}{3}<1$, for all $n\in\N$ we have $0<(\frac{1}{3})^n\le1$. Multiplying by $-7<0$ reverses the inequalities:
\[-7\le-7(\frac{1}{3})^n<0.\]
Adding $8$ and using the previous formula, we obtain $1\le u_n<8$.

\item \textit{[1 pt]} For all $n\in\N$, the explicit formula gives
\[\begin{aligned}u_{n+1}-u_n&=-7(\frac{1}{3})^{n+1}+7(\frac{1}{3})^n\\&=7(1-\frac{1}{3})(\frac{1}{3})^n=\frac{14}{3}(\frac{1}{3})^n>0.\end{aligned}\]
The sequence is therefore strictly increasing.

\item \textit{[0.5 pt]} The claim is false. For each integer $n$, $8-u_n=7(\frac{1}{3})^n>0$. So no term is equal to $8$, whatever the integer $n$.

\end{enumerate}

\subsection*{Exercise 2 · Induction and bounds}\begin{enumerate}

\item \textit{[0.5 pt]} $v_1=\sqrt{42+0}=\sqrt{42}$ and $v_2=\sqrt{42+\sqrt{42}}$.

\item \textit{[1.5 pt]} Let $P(n)$: ``$v_n$ is defined and $0\le v_n\le 7$''.
\par\smallskip\textbf{Base case.} $v_0=0$ is defined and belongs to $[0,7]$; $P(0)$ is true.
\par\smallskip\textbf{Inductive step.} Fix $n\in\N$ and assume $P(n)$ true. Then
\[0\le 42\le 42+v_n\le 42+7=49.\]
The square root is therefore defined and, since this function is increasing on $[0,+\infty[$,
\[0\le v_{n+1}=\sqrt{42+v_n}\le\sqrt{49}=7.\]
Thus $P(n+1)$ is true.
\par\smallskip\textbf{Conclusion.} By induction, all terms are defined and belong to $[0,7]$.

\item \textit{[1.5 pt]} For $n\in\N$, let $R(n)$: ``$v_{n+1}\ge v_n$''.
\par\smallskip\textbf{Base case.} $v_1=\sqrt{42}>0=v_0$; $R(0)$ is true.
\par\smallskip\textbf{Inductive step.} Fix $n\ge0$ and assume $v_{n+1}\ge v_n$. Adding $42$ and then using the fact that the square root is increasing,
\[v_{n+2}=\sqrt{42+v_{n+1}}\ge\sqrt{42+v_n}=v_{n+1}.\]
This is $R(n+1)$.
\par\smallskip\textbf{Conclusion.} By induction, $v_{n+1}\ge v_n$ for all $n\ge0$: $(v_n)$ is increasing.

\item \textit{[1 pt]} Let $S(n)$: ``$v_n<7$''. At index $0$, $v_0=0<7$.
Fix $n\ge0$ and assume $S(n)$ true. The radicands being positive, the strict increase of the square root gives
\[v_{n+1}=\sqrt{42+v_n}<\sqrt{42+7}=\sqrt{49}=7.\]
So $S(n+1)$ is true. The base case and the inductive step establish by induction that $v_n<7$ for all $n\in\N$.

\item \textit{[0.5 pt]} No. For example $P(n):n\ge1$ is true for all integers $n\ge1$ but false at $0$.

\end{enumerate}

\subsection*{Exercise 3 · Vector calculations in a basis}\begin{enumerate}

\item \textit{[1 pt]} In an equality $\vec v=t\vec u$, the coefficient of $\vj $ would be $1=0$, a contradiction by uniqueness of coordinates in a basis.

\item \textit{[0.5 pt]} Expanding, $5\vec u+5\vec v=5(\vi +\vec k)+5(\vj +\vec k)=5\vi +5\vj +10\vec k=\vec w$.

\item \textit{[1.5 pt]} Since $\vec u,\vec v$ are not collinear, coplanarity would be equivalent to $\vec z=\lambda\vec u+\mu\vec v$. Uniqueness of the decomposition in the basis requires
\[\lambda=5,\qquad\mu=5,\qquad\lambda+\mu=11.\]
The first two equalities give $\lambda+\mu=10$, contradicting the third. The three vectors are therefore not coplanar.

\item \textit{[1 pt]} $5\vec u-5\vec v=5(\vi +\vec k)-5(\vj +\vec k)=5\vi -5\vj $.

\item \textit{[1 pt]} $\vect{AB}$ and $\vect{AC}$ are non-collinear and $\vect{AD}$ is their linear combination: $D\in(ABC)$.

\end{enumerate}

\mxneed{24}

\subsection*{Exercise 4 · Geometry of a parallelepiped}\begin{enumerate}

\item \textit{[1 pt]} By Chasles' relation, $\vect{AG}=\vect{AB}+\vect{BC}+\vect{CG}=\vi +\vj +\vec k$. Similarly, $\vect{BH}=\vect{BA}+\vect{AD}+\vect{DH}=-\vi +\vj +\vec k$.

\item \textit{[1 pt]} By Chasles' relation, $\vect{IJ}=\vect{AJ}-\vect{AI}=\frac{\vj -\vi }{15}=\frac{\vect{BD}}{15}$. These vectors are non-zero and collinear: $(IJ)$ and $(BD)$ are parallel.

\item \textit{[1 pt]} We have $\vect{CG}=\vect{AE}=\vec k$ and $\vect{AB}=\vi $. These two basis vectors are not collinear: an equality $\vec k=t\vi $ would contradict the uniqueness of the decomposition in this basis.

\item \textit{[1 pt]} The diagonals of a parallelepiped bisect each other. Thus $\vect{AO}=\frac12\vect{AG}$, which proves $O\in(AG)$.

\item \textit{[1 pt]} The points satisfy $\vect{IJ}=\frac{-\vi +\vj }{15}$ and $\vect{IK}=\frac{-\vi +\vec k}{15}$. If $\vect{IK}=t\vect{IJ}$, the coefficient of $\vec k$ would give $\frac{1}{15}=0$, which is impossible. The vectors are not collinear, so $I,J,K$ are not collinear.

\end{enumerate}

\fi

\fi

